\documentclass{article}
\usepackage[UTF8]{ctex}
\usepackage{amsmath,amssymb,amsfonts}
\usepackage{enumitem}
\usepackage{geometry}
\geometry{a4paper, margin=1in}
% === 错题本排版优化（不参与题目标识） ===
\setlength{\parindent}{0pt}
\setlength{\parskip}{0.35em plus 0.10em minus 0.08em}
\setlength{\emergencystretch}{3em}
\setlist{itemsep=0.25em, topsep=0.40em, parsep=0pt}
\allowdisplaybreaks[2]
% === 错题本排版优化结束 ===
\begin{document}
\section*{湖南大学2026年高等代数真题}
\section*{134．湖南大学2026年高等代数试卷（回忆版）}
（注．题目来自考生回忆，部分表述或数据可能有误）

\subsection{题目1}
1．设

\[
f(x)=\left|\begin{array}{ccccc}
x & & & & 9 \\
-1 & x & & & 33 \\
& -1 & x & & 19 \\
& & -1 & x & 12 \\
& & & -1 & x+6
\end{array}\right| .
\]

将 \(f(x)\) 在有理数域 \(\mathbb{Q}\) 上分解为不可约多项式的乘积。

\subsection{题目2}
2．设 \(f_{k}(x)=1+x+\cdots+x^{k-1}\) ，求

\[
D_{n}=\left|f_{i}(j)\right|=\left|\begin{array}{cccc}
f_{1}(1) & f_{1}(2) & \cdots & f_{1}(n) \\
f_{2}(1) & f_{2}(2) & \cdots & f_{2}(n) \\
\vdots & \vdots & & \vdots \\
f_{n}(1) & f_{n}(2) & \cdots & f_{n}(n)
\end{array}\right| \text {. }
\]

\subsection{题目3}
3．给定有理数域上的线性方程组

\[
\left\{\begin{array}{l}
x_{1}+x_{2}+x_{3}+x_{4}=0 \\
x_{2}+2 x_{3}+2 x_{4}=1 \\
-x_{2}+\lambda x_{3}-2 x_{4}=\mu \\
3 x_{1}+2 x_{2}+x_{3}+(\lambda+3) x_{4}=-1
\end{array}\right.
\]

\begin{itemize}
  \item[（1）] \(\lambda, \mu \in \mathbb{Q}\) 取何值时，线性方程组无解？
  \item[（2）] \(\lambda, \mu \in \mathbb{Q}\) 取何值时，线性方程组有唯一解？并求其解。
  \item[（3）] \(\lambda, \mu \in \mathbb{Q}\) 取何值时，线性方程组有无穷多解？并求其所有解．
\end{itemize}

\subsection{题目4}
4．在复数域上有 \(n\) 阶矩阵

\[
A=\left(\begin{array}{ccccc}
0 & -1 & 0 & \cdots & 0 \\
0 & 0 & -1 & \cdots & 0 \\
\vdots & \vdots & \vdots & & \vdots \\
0 & 0 & 0 & \cdots & -1 \\
-1 & 0 & 0 & \cdots & 0
\end{array}\right) .
\]

求出 \(A\) 的所有复特征值以及其对应的特征子空间。

\subsection{题目5}
5．\(E_{n}\) 是单位矩阵，\(\alpha, \beta \in \mathbb{R}^{n}\) 是列向量，证明：\(E_{n}+\alpha \beta^{\mathrm{T}}\) 可逆的充要条件是 \(1+\beta^{\mathrm{T}} \alpha \neq 0\) ，并求\(E_{n}+\alpha \beta^{\mathrm{T}}\) 的逆．

\subsection{题目6}
6．证明：\(A\) 的行向量组可以由 \(B\) 的行向量组线性表出的充要条件是 \(B X=0\) 的解都是 \(A X=0\) 的解．

\subsection{题目7}
7．已知 \(P\) 是 \(m \times s\) 矩阵，\(Q\) 是 \(t \times n\) 矩阵，\(A\) 是 \(s \times t\) 矩阵，证明：\(r(P A Q) \geq r(P A)+r(A Q)-r(A)\)．

\subsection{题目8}
8．\(A\) 为 \(n\) 阶实对称矩阵，\(p\) 为正整数，\(f(X)=X^{\mathrm{T}} A X\) ，其中 \(X=\left(x_{1}, x_{2}, \cdots, x_{n}\right)^{\mathrm{T}}\) ，证明：\(A\) 的正惯性指数大于等于 \(p\) 的充要条件是存在 \(\mathbb{R}^{n}\) 的 \(p\) 维子空间 \(W\) ，使得 \(f(X)>0, \forall X \in W\) 且 \(X \neq 0\) 。

\subsection{题目9}
9．在欧氏空间中 \(\mathbb{R}^{n}\) ，定义标准内积

\[
(X, Y)=x_{1} y_{1}+x_{2} y_{2}+\cdots+x_{n} y_{n} .
\]

其中 \(X=\left(x_{1}, x_{2}, \cdots, x_{n}\right)^{\mathrm{T}}, Y=\left(y_{1}, y_{2}, \cdots, y_{n}\right)^{\mathrm{T}}, \alpha_{1}, \alpha_{2}, \cdots, \alpha_{m}\) 和 \(\beta_{1}, \beta_{2}, \cdots, \beta_{m}\) 是 \(\mathbb{R}^{n}\) 中的向量组，证明：\(\left\langle\alpha_{i}, \alpha_{j}\right\rangle=\left\langle\beta_{i}, \beta_{j}\right\rangle, 1 \leq i, j \leq m\) 的充要条件是存在正交变换 \(\mathcal{A}\) ，使得 \(\mathcal{A}\left(\alpha_{i}\right)=\beta_{i}, 1 \leq i \leq m\) 。

\subsection{题目10}
10．设 \(M_{n}(\mathbb{C})\) 复数域 \(\mathbb{C}\) 上所有 \(n\) 阶矩阵对通常矩阵加法和数乘运算构成的线性空间，定义 \(M_{n}(\mathbb{C})\)上的线性变换 \(\varphi_{A}(X)=A X A\) ，其中 \(A \in M_{n}(\mathbb{C})\) 。证明：若 \(A\) 可相似对角化，则存在 \(M_{n}(\mathbb{C})\) 的一组基，使得 \(\varphi_{A}\) 在这组基下的矩阵也为对角矩阵。
\end{document}
